% abstract en español
Se presenta el diseño e implementación de un \textit{pipeline} de simulación de
dinámica molecular (MD) para el estudio del rebote inelástico de clústeres de
Lennard-Jones (LJ) sobre una superficie cristalina FCC de argón. A diferencia
de enfoques \textit{ad hoc}, el trabajo enfatiza la arquitectura del software
como componente central de investigación: se emplean metaclases, descriptores,
decoradores paramétricos, generadores, jerarquías de excepciones tipadas y
análisis con \texttt{pandas}/\texttt{scipy}, todo integrado con el motor de MD
\textsc{Lammps}. El observable central es el coeficiente de restitución $e$ en
función del tamaño del clúster $N \in \{1, 2, 7, 13, 55\}$, la velocidad
incidente $v_0$, el ángulo $\theta$ y la adhesión proyectil-superficie
$\varepsilon_{ps}$. Se realizaron \num{1350} simulaciones en barrido paramétrico
paralelo, mostrando que $\alpha_{E,\mathrm{int}}(N)$ crece sistemáticamente con
$N$ mientras el COR decrece, en consistencia con el modelo de disipación por
modos vibracionales internos. Se introduce el número de Weber
$\mathrm{We} = E_{\mathrm{inc}}/(N\varepsilon_{ps})$ como variable de escala
que unifica los regímenes de rebote y adhesión. El código fuente está
completamente versionado y reproducible, constituyendo un caso de estudio
pedagógico en programación científica avanzada.
